%; whizzy paragraph -pdf xpdf -latex ./whizzypdfptex.sh
%; whizzy-paragraph "^\\\\begin{frame}"
% latex beamer presentation.
% platex, latex-beamer でコンパイルすることを想定。 

%     Tokyo Debian Meeting resources
%     Copyright (C) 2009 Junichi Uekawa
%     Copyright (C) 2009 Nobuhiro Iwamatsu

%     This program is free software; you can redistribute it and/or modify
%     it under the terms of the GNU General Public License as published by
%     the Free Software Foundation; either version 2 of the License, or
%     (at your option) any later version.

%     This program is distributed in the hope that it will be useful,
%     but WITHOUT ANY WARRANTY; without even the implied warreanty of
%     MERCHANTABILITY or FITNESS FOR A PARTICULAR PURPOSE.  See the
%     GNU General Public License for more details.

%     You should have received a copy of the GNU General Public License
%     along with this program; if not, write to the Free Software
%     Foundation, Inc., 51 Franklin St, Fifth Floor, Boston, MA  02110-1301 USA

\documentclass[cjk,dvipdfmx,12pt]{beamer}
\usetheme{Tokyo}
\usepackage{monthlypresentation}

%  preview (shell-command (concat "evince " (replace-regexp-in-string
%  "tex$" "pdf"(buffer-file-name)) "&")) 
%  presentation (shell-command (concat "xpdf -fullscreen " (replace-regexp-in-string "tex$" "pdf"(buffer-file-name)) "&"))
%  presentation (shell-command (concat "evince " (replace-regexp-in-string "tex$" "pdf"(buffer-file-name)) "&"))

%http://www.naney.org/diki/dk/hyperref.html
%日本語EUC系環境の時
\AtBeginDvi{\special{pdf:tounicode EUC-UCS2}}
%シフトJIS系環境の時
%\AtBeginDvi{\special{pdf:tounicode 90ms-RKSJ-UCS2}}

\title{第133回東京エリアDebian勉強会 \\　\\Debian GNU/kFreeBSD\\セットアップガイド\\2015年版}
\subtitle{}
\author{Norimitsu Sugimoto (杉本 典充) \\dictoss@live.jp}
\date{2015-11-21}
\logo{\includegraphics[width=8cm]{image200607/openlogo-light.eps}}

\begin{document}

\frame{\titlepage{}}

%\section{}

\begin{frame}{アジェンダ}
  \begin{itemize}
  \item 自己紹介
  \item Debian PortsとDebian GNU/kFreeBSD
  \item Debian GNU/kFreeBSD のインストール
  \item Debian GNU/kFreeBSD 固有の Debian パッケージ
  \item Windows と Debian GNU/kFreeBSD のデュアルブート設定
  \item ハードウェアとソフトウェアのセットアップ
  \item おわりに
  \item 参考資料
  \end{itemize}
\end{frame}

\begin{frame}{自己紹介}
  \begin{itemize}
  \item Norimitsu Sugimoto (杉本 典充)
  \item dictoss@live.jp
  \item Twitter: @dictoss
  \item Debian使って10年以上、FreeBSD使って5年以上
  \item Debian GNU/kFreeBSDが気になっておりウォッチ中
  \item 仕事はソフトウェア開発者をやってます
  \end{itemize}
\end{frame}

%\emtext{}

\begin{frame}[containsverbatim]{Debian PortsとDebian GNU/kFreeBSD}
  \begin{itemize}
  \item Debian Ports
    \begin{itemize}
    \item https://www.debian.org/ports/
    \item 様々なCPUアーキテクチャやカーネルに移植させるプロジェクト
    \item 公式リリースと非公式リリースのものがある
    \end{itemize}
  \item Debian GNU/kFreeBSD
    \begin{itemize}
    \item https://wiki.debian.org/Debian\_GNU/kFreeBSD
    \item FreeBSDカーネルをもつDebianであり、linuxカーネル以外でそれなりに動くものは初めて
    \item Debianでのアーキテクチャ名はkfreebsd-i386,kfreebsd-amd64
    \item 2011年、Debian 6(squeeze)でテクノロジープレビューとしてリリース
    \item 2015年、Debian 8(jessie)ではDropされたため非公式版に出戻りになった
    \end{itemize}
  \end{itemize}
\end{frame}


\emtext{Debian GNU/kFreeBSDのインストール}


\begin{frame}[containsverbatim]{イメージの作成と起動}
  \begin{itemize}
  \item daily-imagesのmini.isoをダウンロードする
  \item ZFSを使う場合はkfreebsd-amd64を選択すること
  \item USBメモリへddする
  \item USBメモリからブートしてインストールする
  \item インストーラの制約
    \begin{itemize}
    \item UEFIブートに対応していない
    \item ディスク形式はMBRのみに対応(GPTは不可)
    \end{itemize}
  \item インストーラは日本語の表示ができないためLANG=Cで進める
  \end{itemize}
\end{frame}

\begin{frame}[containsverbatim]{パーティション構成とファイルシステム}
  \begin{itemize}
  \item パーティション構成には注意
  \item rootパーティションは基本パーティションにしなければならない
  \item 特にwindowsとデュアルブートする場合は残りの基本パーティション数を意識する
  \item kfreebsd-amd64でメモリが多いならZFSでもよい。kfreebsd-i386の利用、またはメモリが少ないならUFSを選択する
  \item UFSを選択するとデフォルトは同期書き込みモードのため、遅く感じるかもしれません
  \end{itemize}
\end{frame}


\emtext{Debian GNU/kFreeBSD固有のDebian パッケージ}


\begin{frame}[containsverbatim]{Debian GNU/kFreeBSD固有の\\Debian パッケージ}
  \begin{itemize}
  \item kfreebsd-image パッケージ
  \item zfsutils パッケージ    
  \item freebsd-utils パッケージ
  \item freebsd-smbfs パッケージ
  \item freebsd-ppp パッケージ
  \item pf パッケージ
  \end{itemize}
\end{frame}


\emtext{Windowsと Debian GNU/kFreeBSDのデュアルブート設定}


\begin{frame}[containsverbatim]{Windowsと Debian GNU/kFreeBSDの\\デュアルブート設定}
  \begin{itemize}
  \item kfreebsdはgrub2を利用します 
  \item \# vi 40\_custom.kfreebsd\\
  menuentry ‘‘Windows (loader)’’ {\\
    insmod part\_msdos\\
    insmod ntfs\\
    set root=(hd0,2)\\
    chainloader +1\\
  }
  \item \# update-grub    
  \end{itemize}
\end{frame}


\emtext{ハードウェアとソフトウェアのセットアップ}


\begin{frame}[containsverbatim]{有線 LAN}
  \begin{itemize}
  \item インタフェース名がNICのベンダーによって変わる(em0,re0)
  \item 設定ファイルはDebian共通で、/etc/network/interfaces
    \begin{itemize}
    \item allow-hotplugはLinuxのudevの機能のため利用できない
    \item auto em0などとしておくと外出先でもDHCPでアドレスを取得しにいくので不便
    \end{itemize}
  \item \# ifup em0={network-name} は利用可能
  \end{itemize}
\end{frame}


\begin{frame}[containsverbatim]{無線 LAN}
  \begin{itemize}
  \item \# vi /etc/apt/sources.list で non-freeを追加    
  \item \#apt-get install kfreebsd-image-10-amd64 firmware-iwlwifi wpasupplicant
  \item 2つのインタフェース
    \begin{itemize}
    \item iwn0 (intel無線LANの場合) : 物理インタフェース
    \item wlan0 : 論理インタフェース \\
      \# ifconfig wlan create wlandev iwn0      
    \end{itemize}
  \item \$ wpa\_passphrase apname1 appassword $>$ wpa\_apname1.conf
  \item \# wpa\_supplicant -i wlan0 -c ./wpa\_apname1.conf
  \item \# dhclient wlan0
  \item つながらないときのポイント
    \begin{itemize}
    \item 2.4GHz帯のアクセスポイントを利用する
    \item \# ifconfig wlan0 -ht40
    \end{itemize}
  \end{itemize}
\end{frame}


\begin{frame}[containsverbatim]{USB モデムによるダイアルアップ}
  \begin{itemize}
  \item \# apt-get install freebsd-ppp usb-modeswitch
  \item \# usb\_modeswitch -c /etc/usb\_modeswitch.d/{config}.conf\\
    を実行することで、USBモデムモードに切り替わる
  \item \# vi /etc/ppp/ppp.conf
  \item \#ppp -foreground {pppname}
  \end{itemize}
\end{frame}


\begin{frame}[containsverbatim]{ビデオドライバ}
  \begin{itemize}
  \item Intel内蔵グラフィック、Radeonシリーズを利用するのがよい
  \item nvidiaはドライバが古いものしかない
  \item \# apt-get install xserver-xorg-video-intel \\
    \# apt-get install xserver-xorg-video-ati
  \item KMSを有効にするとコンソールで解像度が上がり、X Window Systemでも最大限性能を発揮できる
    \begin{itemize}
    \item \# kldunload i915\\
      \# kldload i915kms
    \item \# vi /etc/modules\\
      i915kms
    \end{itemize}
  \end{itemize}
\end{frame}


\begin{frame}[containsverbatim]{locale の再設定}
  \begin{itemize}
  \item \# dpkg-reconfigure locales\\
    日本語の場合は ja\_JP.UTF-8 を選択する
  \end{itemize}
\end{frame}


\begin{frame}[containsverbatim]{X Window System のキーボード設定}
  \begin{itemize}
  \item /etc/X11/xorg.confの設定を行わない場合は英語キーボードとして認識される
  \item FreeBSDにおけるxorgのキーボード検出はhalを利用しているため
  \item 現状はウィンドウマネージャの起動後にウィンドウマネージャのキーボード設定やコマンドラインで対応している
  \end{itemize}
\end{frame}


\begin{frame}[containsverbatim]{web ブラウザ}
  \begin{itemize}
  \item iceweasel パッケージが提供されています
  \item chromium パッケージはkfreebsdにはありません
  \item Flashは本家バイナリはLinux用のみのため、gnashを使う必要がある
  \end{itemize}
\end{frame}


\begin{frame}[containsverbatim]{USB 3.0}
  \begin{itemize}

  \item kfreebsd-image-10.1にはxhci.koがstatic linkしている
  \item 動作は未確認
  \end{itemize}
\end{frame}


\begin{frame}[containsverbatim]{サウンドドライバ}
  \begin{itemize}
  \item FreeBSDはOSSです
  \item 最近のPCの多くはsnd\_hda.koドライバの利用のため概ね鳴ります
  \item pulseaudioは動きます
  \end{itemize}
\end{frame}


\begin{frame}[containsverbatim]{電源関係}
  \begin{itemize}
  \item CPUクロック制御はpowerdで処理される
  \item \$ sysctl dev.cpu.0.freq \\
    dev.cpu.0.freq: 800
  \item バッテリー残量取得は \# /usr/sbin/acpiconf -i 0
  \item サスペンド、ハイバーネートは未確認
  \end{itemize}
\end{frame}


\begin{frame}[containsverbatim]{USB メモリの mount}
  \begin{itemize}
  \item mount\_msdosfsコマンドを使います
  \item 文字コードの指定が必要
  \item \# mount\_msdosfs -L ja\_JP.UTF-8 -D CP932 /dev/da0s1 /mnt/usb
  \item exFAT、NTFSは現状エラーになりmountできない
  \end{itemize}
\end{frame}


\begin{frame}[containsverbatim]{Linux エミュレーション}
  \begin{itemize}
  \item FreeBSDのLinuxバイナリ互換機能が使える
  \item linux-i386のバイナリが実行可能
  \item エミュレートできるLinux-kernel APIの都合のため、squeezeならそこそこ動く
  \item debootstrapでコンテナを作る
  \end{itemize}
\end{frame}


\begin{frame}[containsverbatim]{仮想化}
  \begin{itemize}
  \item 動くと思われるもの
    \begin{itemize}
    \item FreeBSD Jail
    \item qemu
    \end{itemize}
  \item  まだ動かないもの
    \begin{itemize}
    \item VirtualBox
    \item bhyve
    \end{itemize}
  \end{itemize}
\end{frame}


\begin{frame}[containsverbatim]{おわりに}
  \begin{itemize}
  \item Debian GNU/kFreeBSDのセットアップ方法をまとめてみました
  \item OSを勉強するにはいいチャンスだと思います
  \item まずは使ってみるところから初めてみるといいと思います
  \item 使ってみて疑問や動かないところがありましたら、相談に乗ります
  \item がんばってstableリリースを目指しましょう
  \end{itemize}
\end{frame}


\begin{frame}{参考情報}
  \begin{itemize}
  \item 「Debian GNU/kFreeBSD - Debian Wiki」 https://wiki.debian.org/Debian\_GNU/kFreeBSD
  \item 「Debian GNU/kFreeBSD における Jail 構築を試してみた」 http://tokyodebian.alioth.debian.org/pdf/debianmeetingresume201502-presentation-sugimoto.pdf
  \end{itemize}
\end{frame}

\end{document}


;;; Local Variables: ***
;;; outline-regexp: "\\([ 	]*\\\\\\(documentstyle\\|documentclass\\|emtext\\|section\\|begin{frame}\\)\\*?[ 	]*[[{]\\|[]+\\)" ***
;;; End: ***
